\section{引言}
\subsection{研究背景}
在新能源汽车产业快速扩张的背景下，动力电池包已成为整车制造中最具代表性的复杂机电产品之一。它不仅决定车辆的续航、安全和热管理性能，而且在制造环节中直接关联成本、节拍和一致性。与传统零部件装配相比，动力电池包的装配对象更重、工序更多、约束更强，焊接缺陷、绝缘失效、紧固异常或定位偏差等问题都可能在后续使用中演化为可靠性与安全风险 \cite{Kumar2025JMEP,Liang2025JRem}。

近年来，动力电池制造正由单机自动化逐步走向系统级数字化。机器人上料、视觉定位、在线检测、协作工位、数字化追溯与工艺数据闭环已经成为先进电池产线的重要方向 \cite{Ayerbe2021AEM,Park2023RSER,Liu2020Access,Tao2020Sustainability}。然而，现有工程实践中仍存在两个突出问题。其一，装配与检测往往被切分为若干局部环节，导致前一道工序积累的偏差难以及时传递到后续工位，返修与重复检测成本随之上升。其二，虽然机器视觉、深度学习、人机协作、大模型智能体和知识图谱等技术发展迅速，但许多研究仍停留在单点应用层面，缺少围绕具体产品全流程的系统集成路径 \cite{Keshvarparast2024JIM,Ma2024AIR}。

\subsection{问题界定}
如果仅停留在“智能制造技术综述”的层面，便难以说明技术如何真正作用于具体产品。基于此，本文聚焦以下问题：针对动力电池包这一具体对象，装配工艺链条如何展开，数字化手段和关键装备如何嵌入其中，不同智能技术如何形成协同，以及在此基础上如何构建一种可落地的新型装配模式。因此，本文将动力电池包视为一个由结构件、电连接件、热管理部件以及感知与控制部件共同构成的复杂产品系统，并从工艺过程与检测过程出发讨论智能装配问题。

现有研究已经为动力电池制造自动化、人机协作和工业智能化提供了若干可借鉴的方法 \cite{Mayoral2023IJPR,Wang2022RCIM,Gasparetto2021RCIM,Zhang2023JMSYKG}，但仍存在三个需要进一步收束的问题。第一，围绕动力电池包这一具体产品的装配与检测主线，相关论述常被拆散在焊接、拆解、检测或协作制造等不同方向中，缺少统一的产品级分析框架。第二，机器人、机器视觉、知识图谱和大模型智能体等技术的应用边界往往被分开讨论，容易出现“技术很多、系统很弱”的问题。第三，关于新型装配模式的讨论容易停留在概念设想层面，缺少从基本原理到软硬件实现路径的完整论证。针对这些不足，本文将“具体产品、工艺过程、检测节点、关键技术和模式设计”组织为一条连续的论述主线。

\subsection{本文工作与结构安排}
本文的主要工作可以概括为四点。第一，围绕动力电池包这一具体产品，对其装配任务、关键工艺和检测节点进行系统梳理。第二，结合机器人、人机协作、AI、机器视觉、知识图谱和大模型智能体等至少六类技术，分析其在动力电池包装配中的适用场景与协同关系。第三，提出一种知识驱动的人机协同动力电池包智能装配与在线检测模式，并明确其基本原理、层级划分和功能边界。第四，从软硬件系统构成、关键技术路径、应用前景和现实局限四个方面讨论该模式的潜在实施路径与工程意义。全文遵循“问题提出、方案构建、价值分析”的逻辑展开。

全文按照“问题提出到方案落地”的顺序展开。第2节围绕动力电池包装配分析，依次讨论产品构成、装配任务、工艺流程与检测节点；第3节讨论装配过程中的数字化基础与关键智能技术；第4节提出新型装配模式，并说明其系统构成与实现路径；第5节分析该模式的应用前景、主要优势与现实局限；第6节给出结论。

为便于后文展开，图~\ref{fig:intro-placeholder}用于说明动力电池包的结构层级与关键装配对象。图中不仅给出了电池包、模组与箱体之间的层级关系，还展示了 BPMS、BDU、保险丝盒、线束、母排、冷却板和框架等关键部件的位置。由此可以看出，动力电池包装配并不是单一零件的定位问题，而是结构件、电连接件、热管理部件与控制部件的协同集成问题。后文关于工艺节点和检测节点的分析，均建立在这一产品分解基础之上 \cite{Villagrossi2023JRem}。

\begin{figure}[H]
  \centering
  \includegraphics[width=0.76\textwidth]{fig1_battery_pack_components.png}
  \caption{动力电池包关键部件与装配层级关系}
  \label{fig:intro-placeholder}
\end{figure}
